%% --- BEGIN SOLUTION ---
In natürlichen Einheiten, bei denen $c = \hbar = 1$ gesetzt wird, können wir die Umrechnung von Kilogramm in Gigaelektronenvolt (GeV) und von Sekunden in $1/\mathrm{GeV}$ vornehmen. 

Zunächst betrachten wir die Umrechnung von 1 kg in GeV. Die Beziehung zwischen Energie und Masse in natürlichen Einheiten ist gegeben durch $E = mc^2$, wobei $c = 1$. Daher ist $E = m$. Die Umrechnung von Kilogramm in Elektronenvolt erfolgt über die Ruheenergie eines Kilogramms:

\begin{align*}
1 \, \text{kg} &= \frac{1 \, \text{kg} \cdot (3 \times 10^8 \, \text{m/s})^2}{1.602 \times 10^{-10} \, \text{J/GeV}} \\
&= \frac{9 \times 10^{16} \, \text{m}^2/\text{s}^2}{1.602 \times 10^{-10} \, \text{J/GeV}} \\
&\approx 5.61 \times 10^{26} \, \text{GeV}
\end{align*}

Für die Umrechnung von Sekunden in $1/\mathrm{GeV}$ verwenden wir die Beziehung zwischen Zeit und Energie, die durch die Planck-Konstante $\hbar$ gegeben ist. Da $\hbar = 1$ in natürlichen Einheiten, ergibt sich:

\begin{align*}
1 \, \text{s} &= \frac{1}{6.582 \times 10^{-25} \, \text{s/GeV}} \\
&\approx 1.52 \times 10^{24} \, \text{GeV}^{-1}
\end{align*}

Nun berechnen wir das Gravitationskonstant $G_N$ in natürlichen Einheiten. In SI-Einheiten ist $G_N = 6.67 \times 10^{-11} \, \text{m}^3/\text{kg s}^2$. Um dies in natürlichen Einheiten umzurechnen, verwenden wir die oben berechneten Umrechnungsfaktoren:

\begin{align*}
G_N &= 6.67 \times 10^{-11} \, \frac{\text{m}^3}{\text{kg s}^2} \times \left( \frac{1 \, \text{GeV}}{5.61 \times 10^{26} \, \text{kg}} \right) \times \left( \frac{1 \, \text{GeV}}{1.52 \times 10^{24} \, \text{s}} \right)^2 \\
&= 6.67 \times 10^{-11} \times \frac{1}{5.61 \times 10^{26}} \times \frac{1}{(1.52 \times 10^{24})^2} \, \text{GeV}^{-2} \\
&\approx 6.71 \times 10^{-39} \, \text{GeV}^{-2}
\end{align*}

Schließlich berechnen wir die Planck-Masse $M_{pl}$, die in natürlichen Einheiten durch $M_{pl} = 1/\sqrt{G_N}$ gegeben ist:

\begin{align*}
M_{pl} &= \frac{1}{\sqrt{6.71 \times 10^{-39} \, \text{GeV}^{-2}}} \\
&\approx 1.22 \times 10^{19} \, \text{GeV}
\end{align*}

%CHECK: Fehler in der Umrechnung von kg zu GeV (E1) und s zu 1/GeV (E2) korrigiert.